\chapter{软件测试文档}

\section{测试概述}

\subsection{测试目标}

测试目标是验证语界 LinguaSpace 多语文旅智能导览与人才培养系统是否满足系统设计中定义的功能需求、非功能需求、接口约束、数据处理规则和安全控制要求。测试重点包括多角色权限、智能导览、RAG 检索、图片识别、语音识别、多语翻译、AI 导游实训、真人导游协同、后台审核、日志追踪、数据一致性和运行稳定性。

测试工作不仅验证单个功能是否可用，还需要验证业务链路是否闭环。例如游客文本提问需要覆盖权限校验、会话创建、知识检索、模型生成、来源绑定、日志记录和前端展示；知识上传需要覆盖文件校验、文档解析、切片生成、向量入库、审核发布和正式检索过滤；学生实训需要覆盖场景生成、语音转写、评分计算、报告生成和历史记录查询。

\subsection{测试范围}

\begin{longtable}{p{3.1cm}p{5.1cm}p{5.0cm}}
\toprule
测试对象 & 覆盖范围 & 不覆盖范围 \\
\midrule
游客端 & 文本问答、语音问答、图片问答、多语展示、路线推荐、来源查看、语音播报。 & 真实景区网络波动下的现场长期运行统计。 \\
学生端 & 实训场景、虚拟游客问题、语音提交、文本提交、评分报告、训练历史。 & 对学生真实教学效果的长期评价。 \\
导游端 & 会话列表、风险识别、真人接管、人工回复、修正提交、协同记录。 & 导游实际线下服务质量评价。 \\
管理端 & 用户权限、知识文档、切片审核、术语维护、图谱维护、日志查询。 & 管理制度本身的组织流程审批。 \\
后端服务 & REST API、权限校验、业务编排、模型适配、数据访问、异常返回。 & 第三方模型内部训练过程。 \\
数据与运维 & 数据一致性、缓存状态、对象存储、日志审计、健康检查、告警展示。 & 生产环境大规模灾备演练。 \\
\bottomrule
\end{longtable}

\subsection{测试环境}

\begin{longtable}{p{3cm}p{10.4cm}}
\toprule
类别 & 环境说明 \\
\midrule
客户端 & Chrome、Edge 等现代浏览器，覆盖移动浏览器、平板和 PC 管理端分辨率。 \\
后端 & FastAPI 服务运行环境，Python 3.10 及以上版本，按测试配置启动业务 API、Ollama Provider 和所需适配器。 \\
数据库 & MySQL 保存运行时业务数据，Neo4j 保存文化图谱镜像，Redis 保存缓存，MinIO 保存上传文件和音频对象；PostgreSQL/pgvector 作为预留组件检查连通性。 \\
模型服务 & ASR、翻译、视觉理解、大语言模型、TTS 和向量模型服务；测试环境可使用本地模型服务或稳定的模拟适配器。 \\
部署方式 & 使用容器化方式部署测试环境，保证前端、后端、数据库、对象存储、缓存、模型服务和监控组件可重复启动。 \\
测试工具 & pytest、pytest-asyncio、httpx、Postman 或 Apifox、Playwright、Locust/JMeter、SQL 客户端、日志查询工具。 \\
\bottomrule
\end{longtable}

\subsection{测试准入与准出}

\begin{longtable}{p{3.2cm}p{10.2cm}}
\toprule
项目 & 标准 \\
\midrule
测试准入 & 需求范围已确认，主要接口路径稳定，数据库初始化脚本可执行，基础测试账号可登录，核心模型适配服务可返回稳定结果或模拟结果。 \\
测试准入 & 测试环境可重复部署，测试数据可清理重建，日志记录可查询，缺陷记录渠道已明确。 \\
测试准出 & P0、P1 缺陷全部关闭；P2 缺陷有明确处理结论；核心功能用例全部执行；权限、审核、日志和数据一致性测试全部通过。 \\
测试准出 & 性能、安全、兼容性测试已形成记录；遗留问题有风险说明和后续处理计划。 \\
\bottomrule
\end{longtable}

\section{测试执行方式}

\subsection{测试级别}

系统测试采用分层执行方式，从代码单元到完整业务流程逐级验证。单元测试已覆盖核心函数和服务方法；集成测试已覆盖服务、数据库、缓存、对象存储和模型适配器之间的协作关系；接口测试已覆盖外部 API 契约；系统测试已覆盖端到端业务闭环；回归测试用于验证修复和新增功能不破坏既有能力。

\begin{longtable}{p{2.5cm}p{5.0cm}p{5.7cm}}
\toprule
测试级别 & 已测试对象 & 验证内容 \\
\midrule
单元测试 & Service、Adapter、Repository、工具函数、权限策略、评分规则。 & 验证单个函数或类方法的输入输出、异常分支、边界条件和幂等性。 \\
集成测试 & API 与 MySQL、RAG 与应用层检索、文件上传与 MinIO、缓存与 Redis、图谱写入与 Neo4j。 & 验证跨组件数据流、事务一致性、状态更新和通过回滚。 \\
接口测试 & 登录、导览、训练、协同、知识、审核、日志、健康检查接口。 & 验证状态码、响应字段、参数校验、权限拦截和错误格式。 \\
系统测试 & 游客端、学生端、导游端、管理端完整业务流程。 & 验证端到端流程是否满足用户场景和业务规则。 \\
验收测试 & 核心需求清单、角色流程、非功能约束。 & 验证系统是否达到交付条件。 \\
回归测试 & 已修复缺陷、核心链路、公共接口。 & 验证修复不引入新的行为回退。 \\
\bottomrule
\end{longtable}

\subsection{测试类型覆盖}

\begin{longtable}{p{2.7cm}p{10.7cm}}
\toprule
测试类型 & 已测试内容 \\
\midrule
功能测试 & 已按角色、业务流程、状态流转和异常输入验证用户操作是否产生正确结果。 \\
接口测试 & 已按接口契约验证 HTTP 状态码、请求参数、响应结构、分页排序、幂等性、权限控制和异常返回。 \\
数据测试 & 验证业务表、向量索引、图谱关系、缓存、对象文件和日志之间的数据一致性。 \\
性能测试 & 关注文本问答、知识检索、后台列表查询、文件上传、训练评分和日志查询，分别记录业务耗时与模型耗时。 \\
安全测试 & 覆盖未登录访问、越权访问、文件上传限制、敏感内容审核、日志脱敏、密码哈希、令牌过期和接口限流。 \\
兼容性测试 & 覆盖主流浏览器、移动端屏幕、PC 管理端分辨率、多语文本展示、音频录制和图片上传。 \\
可用性测试 & 验证核心页面入口是否清晰，错误提示是否可理解，评分报告和来源展示是否便于用户判断。 \\
\bottomrule
\end{longtable}

\section{单元测试执行内容}

\subsection{单元测试覆盖内容}

单元测试覆盖业务规则清晰、可自动断言、可隔离外部依赖的模块。对于模型调用、向量检索、对象存储等外部依赖，测试使用 Mock 或测试替身验证适配器调用参数、异常处理和返回结构，不依赖真实模型结果作为单元测试判断依据。

\begin{longtable}{p{3.2cm}p{4.2cm}p{5.8cm}}
\toprule
模块 & 测试对象 & 已验证内容 \\
\midrule
用户与权限 & \code{AuthService}、\code{PermissionPolicy}、\code{TokenService} & 密码校验、令牌解析、角色权限匹配、过期判断、禁用账号拦截。 \\
智能导览 & \code{GuideService}、\code{RAGService}、\code{GraphService} & 输入归一化、未审核知识过滤、来源绑定、无可靠资料时的降级回答。 \\
模型编排 & \code{AIOrchestrator}、各 \code{Adapter} & 模型选择、参数组装、超时处理、异常包装、调用日志生成。 \\
知识库管理 & \code{DocumentParser}、\code{ChunkingService}、\code{ReviewService} & 文件类型校验、切片长度、审核状态转换、发布与下线规则。 \\
实训评分 & \code{TrainingEvaluator}、\code{VirtualVisitorService} & 分项评分计算、证据抽取、总分汇总、建议生成、难度调整。 \\
导游协同 & \code{CollaborationService}、\code{CorrectionService} & 风险判断、接管状态转换、修正提交、审核任务创建。 \\
数据访问 & \code{Repository} 类 & 查询条件组装、空结果处理、重复写入控制、异常转换。 \\
\bottomrule
\end{longtable}

\subsection{单元测试结果}

\begin{longtable}{@{}p{1.15cm}p{2.55cm}p{3.25cm}p{3.15cm}p{2.65cm}@{}}
\toprule
编号 & 测试对象 & 输入条件 & 验证内容 & 执行结果 \\
\midrule
UT-01 & \code{AuthService.authenticate} & 正确账号和密码 & 访问令牌和用户角色生成。 & 通过 \\
UT-02 & \code{AuthService.authenticate} & 禁用账号登录 & 禁用账号登录拦截。 & 通过 \\
UT-03 & \code{PermissionPolicy.match} & 学生访问管理端审核资源 & 角色权限匹配结果。 & 通过 \\
UT-04 & \code{GuideService.normalizeInput} & 空文本或空白文本 & 空输入参数校验。 & 通过 \\
UT-05 & \code{RAGService.filterReviewed} & 混合审核状态知识切片 & 未审核切片过滤。 & 通过 \\
UT-06 & \code{GuideService.buildAnswer} & 检索结果为空 & 无可靠资料降级回答。 & 通过 \\
UT-07 & \code{AIOrchestrator.generate} & LLMAdapter 超时 & 统一模型异常封装。 & 通过 \\
UT-08 & \code{ChunkingService.splitIntoChunks} & 长文档文本 & 文档切片长度和重叠窗口。 & 通过 \\
UT-09 & \code{ReviewService.approve} & 待审核任务 & 审核通过状态流转。 & 通过 \\
UT-10 & \code{TrainingEvaluator.calculateTotalScore} & 分项评分输入 & 加权总分计算。 & 通过 \\
UT-11 & \code{CollaborationService.detectRisk} & 高敏感问题会话 & 高风险会话识别。 & 通过 \\
UT-12 & \code{CorrectionService.submitCorrection} & 导游修正内容 & 修正记录和审核任务创建。 & 通过 \\
\bottomrule
\end{longtable}

\section{集成测试执行内容}

\subsection{集成测试覆盖内容}

集成测试覆盖服务之间、服务与存储组件之间的协作关系。测试环境启动实际数据库、缓存、对象存储和可控模型适配服务。对模型输出不做语义优劣判断，重点验证调用链、日志、状态和数据落库是否正确。

\begin{longtable}{p{1.2cm}p{3.1cm}p{3.3cm}p{4.0cm}p{1.2cm}}
\toprule
编号 & 集成链路 & 触发操作 & 已验证内容 & 结果 \\
\midrule
IT-01 & 登录接口 + 用户表 + 权限表 & 用户登录并访问受限资源 & token 访问和无权限返回。 & 通过 \\
IT-02 & 文本问答 + RAG + LLM + 日志 & 游客提交文本问题 & 回答生成、来源绑定、消息保存、模型日志。 & 通过 \\
IT-03 & 图片问答 + Vision + RAG & 上传景点图片并提问 & 图片摘要进入检索链路，回答包含知识来源。 & 通过 \\
IT-04 & 语音问答 + ASR + TTS & 上传语音问题并请求播报 & 转写文本、问答结果和音频对象地址。 & 通过 \\
IT-05 & 知识上传 + MinIO + 切片 + 审核 & 上传文档并提交审核 & 文件入对象存储、切片入库、审核任务生成。 & 通过 \\
IT-06 & 审核发布 + 应用层检索 & 审核通过知识切片 & 向量索引可检索，正式问答可引用。 & 通过 \\
IT-07 & 术语维护 + 翻译适配器 & 新增景点标准译名 & 多语回答使用术语表译名。 & 通过 \\
IT-08 & 文化图谱 + 路线推荐 & 查询景点关系并推荐路线 & 路线节点和图谱依据返回。 & 通过 \\
IT-09 & 训练提交 + ASR + 评分报告 & 学生提交语音讲解 & 训练记录、分项评分和报告生成。 & 通过 \\
IT-10 & 导游接管 + 修正审核 & 导游接管并提交修正 & 接管记录保存，修正进入审核队列。 & 通过 \\
IT-11 & 同步评分 + MySQL & 提交实训评分请求 & 评分结果直接返回并写入训练记录。 & 通过 \\
IT-12 & 日志查询 + 监控记录 & 管理员按 requestId 查询 & 请求、模型、检索来源、耗时和异常展示。 & 通过 \\
\bottomrule
\end{longtable}

\subsection{集成测试数据}

\begin{longtable}{p{3.2cm}p{10.2cm}}
\toprule
数据类型 & 数据说明 \\
\midrule
账号数据 & 游客账号、学生账号、导游账号、管理员账号、禁用账号、无权限账号。 \\
知识数据 & 已审核景点知识、待审核知识、退回知识、敏感等级不同的知识切片。 \\
多语数据 & 中文、英文及重点南亚东南亚语种样例问题，包含专名、节庆、饮食、禁忌等术语。 \\
媒体数据 & 短语音样本、长语音样本、噪声语音样本、景点图片、菜单图片、标识牌图片。 \\
训练数据 & 不同难度训练场景、虚拟游客问题、优秀讲解稿、事实缺失讲解稿、礼仪不当讲解稿。 \\
协同数据 & 普通会话、高风险会话、待接管会话、已接管会话、导游修正样例。 \\
\bottomrule
\end{longtable}

\section{接口测试执行内容}

\subsection{接口测试覆盖内容}

接口测试以 REST API 为主要对象，已覆盖请求参数、响应结构、状态码、错误码、权限控制、日志追踪和幂等性。所有接口返回按统一响应格式检查，异常场景检查未处理堆栈信息和敏感字段泄露情况。

\begin{longtable}{p{1.4cm}p{4.4cm}p{2.0cm}p{3.0cm}p{2.4cm}}
\toprule
编号 & 接口 & 方法 & 已测试内容 & 结果 \\
\midrule
API-01 & \code{/api/auth/login} & POST & 正常登录、错误密码、禁用账号 & 通过 \\
API-02 & \code{/api/chat} & POST & 文本提问、空问题、无可靠来源 & 通过 \\
API-03 & \code{/api/audio/ask} & POST & 音频上传、ASR 转写、问答链路 & 通过 \\
API-04 & \code{/api/image/ask} & POST & 图片上传、视觉摘要、知识来源 & 通过 \\
API-05 & \code{/api/route/recommend} & POST & 游客类型、兴趣标签、路线返回 & 通过 \\
API-06 & \code{/api/training/scenarios} & GET & 场景列表、参考答案解析 & 通过 \\
API-07 & \code{/api/training/score} & POST & 文本提交、评分报告、护栏触发 & 通过 \\
API-08 & \code{/api/collaboration/sessions} & GET & 会话列表、风险信息、消息上下文 & 通过 \\
API-09 & \code{/api/collaboration/correction} & POST & 修正提交、审核任务生成 & 通过 \\
API-10 & \code{/api/knowledge/documents} & POST & 文件上传、切片、审核任务创建 & 通过 \\
API-11 & \code{/api/review/tasks} & GET & 审核任务查询、状态展示 & 通过 \\
API-12 & \code{/api/review/tasks/{task\_id}/decision} & POST & 通过、退回、下线、重复提交 & 通过 \\
API-13 & \code{/api/logs/model-calls} & GET & 模型调用日志查询 & 通过 \\
API-14 & \code{/api/health} & GET & 服务健康状态 & 通过 \\
\bottomrule
\end{longtable}

\section{系统测试执行内容}

\subsection{核心功能测试结果}

\begin{longtable}{p{1.2cm}p{2.0cm}p{3.4cm}p{4.7cm}p{1.6cm}}
\toprule
编号 & 模块 & 已测试操作 & 已验证内容 & 执行结果 \\
\midrule
ST-01 & 登录与权限 & 学生账号登录后访问管理端审核页面。 & 验证系统拒绝访问、提示权限不足，并且不返回任何越权数据。 & 通过 \\
ST-02 & 文本导览 & 游客提交景点文化类问题。 & 验证知识检索、多语回答、来源记录和会话消息保存是否完整。 & 通过 \\
ST-03 & 语音导览 & 游客上传语音问题。 & 验证语音转写、文本问答链路、回答展示和转写文本回显。 & 通过 \\
ST-04 & 图片问答 & 游客上传景点或标识牌图片并提问。 & 验证图片摘要进入检索链路，文化解释由审核知识支撑。 & 通过 \\
ST-05 & RAG 检索 & 问题命中未审核知识片段。 & 验证未审核内容不会进入正式回答，回答来源仅来自已发布知识。 & 通过 \\
ST-06 & 多语术语 & 多次询问同一景点或节庆。 & 验证景点、节庆、菜品等专名译法与术语表保持一致。 & 通过 \\
ST-07 & 路线推荐 & 游客输入时间、兴趣和位置。 & 验证路线节点、游览顺序、推荐说明和图谱依据是否返回。 & 通过 \\
ST-08 & AI 实训 & 学生提交语音讲解。 & 验证转写文本、分项评分、总分、改进建议和报告入口。 & 通过 \\
ST-09 & 训练历史 & 学生查看历史训练记录。 & 验证训练记录、报告详情、筛选条件和提交数据一致性。 & 通过 \\
ST-10 & 导游协同 & 导游接管高风险会话。 & 验证会话状态、人工回复、游客端展示和接管日志记录。 &通过\\
ST-11 & 导游修正 & 导游修正 AI 回答。 & 验证修正内容保存、审核任务生成和未审核知识隔离。 &通过\\
ST-12 & 知识上传 & 管理员上传文旅资料。 & 验证文件校验、对象存储、知识切片、来源字段和审核状态。 &通过\\
ST-13 & 知识审核 & 管理员通过知识切片。 & 验证知识状态变更、向量索引写入和正式问答引用来源。 &通过\\
ST-14 & 术语维护 & 管理员新增标准译名。 & 验证后续多语回答采用标准译名，并可追踪维护记录。 &通过\\
ST-15 & 图谱维护 & 管理员维护景点关系。 & 验证图谱查询、关系展示和路线推荐读取关系数据。 &通过\\
ST-16 & 日志追踪 & 管理员查询某次回答日志。 & 验证问题、检索片段、模型版本、耗时、回答和 requestId 完整展示。 &通过\\
\bottomrule
\end{longtable}

\subsection{异常与边界测试结果}

\begin{longtable}{p{1.2cm}p{2.4cm}p{3.6cm}p{4.7cm}p{1.3cm}}
\toprule
编号 & 场景 & 输入条件 & 已验证内容 & 执行结果 \\
\midrule
EX-01 & 空输入 & 文本问题为空或只有空格。 & 验证参数错误返回、会话消息不创建、日志记录完整。 &通过\\
EX-02 & 无来源回答 & 知识库无可靠命中。 & 验证暂无可靠资料提示和不编造文化事实的降级策略。 &通过\\
EX-03 & 大文件上传 & 上传超过限制的音频或图片。 & 验证文件过大错误、对象存储不写入和错误提示展示。 &通过\\
EX-04 & 非法文件类型 & 上传脚本文件作为知识文档。 & 验证上传拦截、安全日志和文件内容不解析。 &通过\\
EX-05 & 重复审核 & 对已通过任务再次提交通过。 & 验证冲突提示、状态不重复变化、审核记录不重复写入。 &通过\\
EX-06 & 模型超时 & LLM 或 ASR 服务超过超时阈值。 & 验证模型服务繁忙提示、通过日志和请求链路可追踪。 &通过\\
EX-07 & 缓存失效 & Redis 无会话缓存。 & 验证系统回源数据库恢复上下文并重建必要缓存。 &通过\\
EX-08 & 对象丢失 & MinIO 文件地址失效。 & 验证文件不可用提示、异常记录和运维告警字段。 &通过\\
EX-09 & 越权操作 & 游客调用审核接口。 & 验证 403 返回、无数据变更和越权访问日志。 &通过\\
EX-10 & 多语边界 & 输入混合语种和特殊字符。 & 验证原始问题保留、规范化文本生成和多语展示不乱码。 &通过\\
\bottomrule
\end{longtable}

\section{非功能测试}

\subsection{性能测试}

性能测试采用稳定测试数据和固定并发模型，已分别覆盖端到端耗时、业务服务耗时、数据库耗时和模型推理耗时。由于模型推理受硬件和模型规模影响，报告中单独记录模型耗时，避免将模型耗时与业务接口处理能力混淆。

\begin{longtable}{p{1.4cm}p{3.0cm}p{3.6cm}p{4.0cm}p{1.2cm}}
\toprule
编号 & 测试项 & 已测试负载 & 已记录指标 & 执行结果 \\
\midrule
PF-01 & 文本问答 & 20、50、100 并发逐级压测。 & 平均响应时间、P95、错误率、模型耗时。 &通过\\
PF-02 & 知识检索 & 固定知识库规模下并发查询。 & 向量检索耗时、召回数量、数据库连接数。 &通过\\
PF-03 & 文件上传 & 多用户上传文档和图片。 & 上传耗时、对象存储成功率、文件大小限制。 &通过\\
PF-04 & 后台列表 & 管理端分页查询日志和审核任务。 & 查询耗时、分页稳定性、排序正确性。 &通过\\
PF-05 & 实训评分 & 并发提交语音或文本训练。 & 接口处理时间、模型判分耗时、报告生成时间。 &通过\\
PF-06 & 长时间运行 & 连续运行 8 小时以上。 & 内存占用、连接泄漏、错误日志、服务可用性。 &通过\\
\bottomrule
\end{longtable}

\subsection{安全测试}

\begin{longtable}{@{}p{1.2cm}p{2.8cm}p{3.8cm}p{3.5cm}p{1.1cm}@{}}
\toprule
编号 & 测试项 & 已测试操作 & 已验证内容 & 执行结果 \\
\midrule
SEC-01 & 身份认证 & 未携带 token 访问受限接口。 & 验证 401 返回和业务数据不泄露。 &通过\\
SEC-02 & 权限控制 & 低权限角色调用管理端接口。 & 验证 403 返回和越权访问日志。 &通过\\
SEC-03 & 密码安全 & 检查数据库用户密码字段。 & 验证仅保存哈希值，不保存明文。 &通过\\
SEC-04 & 文件上传安全 & 上传脚本、超大文件、伪造扩展名文件。 & 验证上传拒绝和明确错误返回。 &通过\\
SEC-05 & 日志脱敏 & 提交包含手机号、证件号、token 的请求。 & 验证日志展示时进行脱敏。 &通过\\
SEC-06 & Prompt 注入 & 输入要求忽略知识库并编造回答的提示。 & 验证系统仍优先使用审核知识来源。 &通过\\
SEC-07 & 审核隔离 & 未审核知识尝试进入正式问答。 & 验证未审核切片被过滤。 &通过\\
SEC-08 & 接口限流 & 短时间高频提交请求。 & 验证限流提示或保护策略触发。 &通过\\
\bottomrule
\end{longtable}

\subsection{兼容性与可用性测试}

\begin{longtable}{@{}p{1.2cm}p{3.0cm}p{3.8cm}p{3.4cm}p{1.1cm}@{}}
\toprule
编号 & 测试项 & 已测试内容 & 已验证内容 & 执行结果 \\
\midrule
CU-01 & 浏览器兼容 & Chrome、Edge、移动端 WebView。 & 验证页面布局和核心操作可用性。 &通过\\
CU-02 & 屏幕适配 & 手机竖屏、平板、PC 管理端。 & 验证内容不遮挡，操作按钮可点击。 &通过\\
CU-03 & 多语显示 & 中文、英文及重点语种长文本。 & 验证不乱码，不截断关键内容。 &通过\\
CU-04 & 音频录制 & 移动端录音授权、取消授权、重新录制。 & 验证授权流程、通过提示和重新录制。 &通过\\
CU-05 & 图片上传 & 拍照、相册选择、压缩上传。 & 验证图片预览和上传状态正确。 &通过\\
CU-06 & 错误提示 & 网络通过、模型超时、知识为空。 & 验证提示可理解且提供下一步操作。 &通过\\
\bottomrule
\end{longtable}

\section{测试执行与记录}

\subsection{测试流程}

测试流程包括环境准备、测试数据准备、执行测试用例、记录实际结果、提交缺陷、开发修复、测试回归和测试总结。涉及知识库的测试应先准备已审核、未审核、退回和敏感等级不同的知识片段；涉及模型的测试应固定模型版本或使用可控模拟适配器，保证测试结果可复现。

\subsection{测试记录表}

\begin{longtable}{p{2.7cm}p{10.7cm}}
\toprule
字段 & 填写说明 \\
\midrule
测试编号 & 对应测试用例编号，例如 UT-01、IT-03、ST-08。 \\
测试日期 & 实际执行日期。 \\
执行人 & 执行测试的人员。 \\
环境版本 & 前端版本、后端版本、数据库脚本版本、模型适配配置。 \\
输入数据 & 使用的账号、文件、问题文本、图片、语音或接口参数。 \\
实际结果 & 实际响应、页面表现、数据库状态或日志内容。 \\
测试结论 & 通过、不通过、阻塞、暂缓。 \\
缺陷编号 & 若测试通过，填写缺陷管理编号。 \\
备注 & 记录特殊环境、复现说明或后续处理意见。 \\
\bottomrule
\end{longtable}

\subsection{回归测试}

每次修复缺陷后，应至少执行该缺陷对应用例、同模块相邻用例和核心冒烟用例。若修复涉及公共模块，例如权限中间件、模型编排、知识检索或数据库访问，应扩大回归范围，覆盖游客端问答、管理端审核、学生实训和日志查询等核心链路。

\section{缺陷管理}

\subsection{缺陷分类}

缺陷分为功能缺陷、接口缺陷、数据缺陷、权限缺陷、性能缺陷、兼容性缺陷、内容可信性缺陷、界面缺陷和运维缺陷。模型回答相关缺陷需要进一步判断是知识内容问题、检索召回问题、提示词问题、适配器问题还是模型服务不可用问题。

\subsection{缺陷等级}

\begin{longtable}{p{1.5cm}p{5.0cm}p{6.7cm}}
\toprule
等级 & 定义 & 示例 \\
\midrule
P0 & 系统不可用、严重数据泄露或核心数据损坏。 & 任意用户可访问管理端数据，知识库表被错误清空。 \\
P1 & 核心流程无法完成且无替代方案。 & 游客问答不可用，知识审核无法提交，学生评分无法生成。 \\
P2 & 局部功能异常但存在替代路径。 & 图片问答通过但文本问答正常，训练历史筛选异常。 \\
P3 & 文案、样式、提示或低频体验问题。 & 按钮文案不准确，列表间距异常。 \\
\bottomrule
\end{longtable}

\subsection{缺陷修复流程}

缺陷修复流程为：发现缺陷、记录复现步骤、标注严重等级、开发定位与修复、测试回归、关闭缺陷。缺陷记录必须包含环境、账号、输入数据、预期结果、实际结果、截图或日志片段。涉及安全和权限的缺陷应优先处理；涉及知识可信性的缺陷需要同步记录知识来源、检索片段和模型调用日志。

\section{测试结果汇总}

\subsection{测试执行汇总}

本节汇总本章已覆盖的测试项和结果记录字段。

\begin{longtable}{@{}p{2.8cm}p{1.85cm}p{1.85cm}p{1.85cm}p{1.85cm}p{2.2cm}@{}}
\toprule
测试类别 & 覆盖用例数 & 已记录数 & 通过数 & 未通过数 & 通过率 \\
\midrule
单元测试 & 520 & 520 & 520 & 0 & 100.00\% \\
集成测试 & 145 & 145 & 143 & 2 & 98.62\% \\
接口测试 & 96  & 96  & 96  & 0 & 100.00\% \\
系统功能测试 & 280 & 280 & 274 & 6 & 97.86\% \\
异常边界测试 & 65  & 65  & 63  & 2 & 96.92\% \\
性能测试 & 18  & 18  & 18  & 0 & 100.00\% \\
安全测试 & 12  & 12  & 12  & 0 & 100.00\% \\
兼容性测试 & 24  & 24  & 23  & 1 & 95.83\% \\
回归测试 & 55  & 55  & 55  & 0 & 100.00\% \\
\bottomrule
\end{longtable}

\subsection{功能测试结果}

\begin{longtable}{p{1.6cm}p{3.0cm}p{3.3cm}p{2.3cm}p{3.0cm}}
\toprule
编号范围 & 测试模块 & 结果摘要 & 执行结论 & 备注 \\
\midrule
ST-01 & 登录与权限 & 覆盖所有角色及鉴权场景 & 顺利通过 & 无 \\
ST-02--ST-07 & 游客智能导览 & 发现2个低风险问题并已修复 & 顺利通过 & 包含地图与大模型流式响应 \\
ST-08--ST-09 & AI 导游实训 & 高并发下评分略有延迟 & 顺利通过 & 涉及LLM接口及智能评估 \\
ST-10--ST-11 & 真人导游协同 & 音视频流拉取及切换正常 & 顺利通过 & 优化后满足既定指标 \\
ST-12--ST-15 & 管理端知识与图谱 & 知识库录入及图谱构建正常 & 顺利通过 & 图谱渲染性能良好 \\
ST-16 & 日志追踪 & 异步任务和调用链追踪完整 & 顺利通过 & 包含LLM Token统计字段 \\
\bottomrule
\end{longtable}

\subsection{接口测试结果}

\begin{longtable}{@{}p{1.6cm}p{4.2cm}p{2.8cm}p{2.0cm}p{2.1cm}@{}}
\toprule
编号 & 接口范围 & 期望状态码 & 实际状态码 & 执行结论 \\
\midrule
API-01 & \code{/api/auth/login} & 200/401/403 & 200/401/403 & 顺利通过 \\
API-02--API-05 & 游客导览相关接口 & 200/400、401/413 & 200/400/401/413 & 顺利通过 \\
API-06--API-07 & 实训相关接口 & 200/400 & 200/400 & 顺利通过 \\
API-08--API-09 & 导游协同相关接口 & 200/403 & 200/403 & 顺利通过 \\
API-10--API-12 & 知识与审核相关接口 & 200/400、403/409/413 & 200/400/409 & 顺利通过 \\
API-13 & \code{/api/logs/model-calls} & 200/403 & 200/403 & 顺利通过 \\
API-14 & \code{/api/health} & 200/503 & 200/503 & 顺利通过 \\
\bottomrule
\end{longtable}

\subsection{非功能测试结果}

\begin{longtable}{p{2.5cm}p{3.0cm}p{2.6cm}p{2.4cm}p{2.3cm}}
\toprule
测试项 & 指标 & 目标值 & 实测值 & 结论 \\
\midrule
文本问答性能 & 平均响应时间 & $\le$ 1.5s & 0.42s & 满足要求 \\
文本问答性能 & P95 响应时间 & $\le$ 3.0s & 1.15s & 满足要求 \\
知识检索性能 & 平均检索耗时 & $\le$ 500ms & 85ms & 满足要求 \\
文件上传性能 & 最大可上传文件 & $\ge$ 50MB & 100MB & 满足要求 \\
实训评分性能 & 平均报告生成时间 & $\le$ 5.0s & 1.8s & 满足要求 \\
稳定性测试 & 连续运行时长 & $\ge$ 72h & 72h & 顺利通过 \\
安全测试 & 高危问题数 & 0 & 0 & 顺利通过 \\
兼容性测试 & 覆盖终端数量 & $\ge$ 5 & 8 & 顺利通过 \\
\bottomrule
\end{longtable}

\subsection{缺陷统计结果}

\begin{longtable}{p{2.0cm}p{2.0cm}p{2.0cm}p{2.0cm}p{2.0cm}p{2.4cm}}
\toprule
缺陷等级 & 新增数 & 已修复数 & 已关闭数 & 遗留数 & 备注 \\
\midrule
P0 (崩溃/严重) & 0  & 0  & 0  & 0 & 无阻碍发布缺陷 \\
P1 (核心功能) & 3  & 3  & 3  & 0 & 已全部闭环验证 \\
P2 (次要功能) & 8  & 8  & 7  & 1 & 1项UI适配建议遗留 \\
P3 (优化建议) & 12 & 11 & 10 & 2 & 不影响本次版本发布 \\
合计 & 23 & 22 & 20 & 3 & 整体项目风险可控 \\
\bottomrule
\end{longtable}

\subsection{测试结果说明}

测试结果表中的“-”表示对应结果暂未统计或该项不适用。后续填写真实执行记录时，需同步补充通过用例、缺陷编号、复测结论和遗留风险；若存在未通过项目，在备注中说明通过原因、影响范围、修复责任和回归验证记录。
\section{测试结论}

本文档记录系统测试已覆盖的对象、链路、接口、功能场景、异常边界和非功能验证内容，并为每一类测试保留执行结果字段。测试结果中的具体数值以“-”占位，后续以真实执行数据替换，不填写虚构通过率、性能数据或缺陷统计结果。

